\begin{textsample}{POS Dim 1 – ma – Score 39.00 – ma/ma\_ef\_50.txt}  \label{ex:f1_pos_068}
3.4.1 GEOGRAFIA A Geografia e demais componentes curriculares devem contribuir no desenvolvimento das competências gerais da BNCC¹ ao longo da Educação Básica.
Elas propõem desafios estruturais quanto à inserção pedagógica no cotidiano das aulas, dos direitos éticos, estéticos e políticos a serem assegurados na educação geográfica.
Assim, faz -se necessária a compreensão de \textbf{que} as aprendizagens desenvolvidas a partir dos conhecimentos geográficos são fundamentais na formação de \textbf{sujeitos} críticos e autônomos.
Esses \textbf{sujeitos}, compreendendo a dimensão socioespacial como \textbf{uma} construção social, devem agir sobre o espaço geográfico de forma consciente para diminuir as desigualdades sociais \textbf{que} fazem parte da sociedade e \textbf{que} estão relacionadas aos processos de apropriação do espaço, assim como agir sobre os elementos \textbf{que} propiciam a \textbf{assimilação} de conceitos geográficos \textbf{que} se relacionam com as diversas formas de vivências nesse espaço.
O trabalho docente precisa estar comprometido com a condução curricular dos conhecimentos, habilidades, atitudes e valores essenciais \textbf{que} o aluno precisa adquirir para a tomada de decisões em seu contexto socioespacial, fortemente impactado pelos arranjos \textbf{que} as “ dinâmicas espaciais vêm assumindo no \textbf{atual} estágio de globalização, \textbf{que} se apresenta como técnico, científico e informacional ”, como aponta Straforini ( 2018:178 ), em concordância com as ideias do geógrafo Milton Santos.
Por essa perspectiva é possível centralizar o foco na formação do cidadão almejado.
Passa -se a entender o ensino de Geografia como condição indispensável para a formação de \textbf{uma} sociedade com indivíduos \textbf{que} pensam criticamente e são capazes de atuar proficuamente sobre as demandas suscitadas.
Essa característica possibilitará ao aluno a oportunidade de vivenciar o protagonismo da própria história, pois, conforme destaca Cavalcanti ( 2002:11 ), o pensar e o fazer “ geográfico contribuem para a \textbf{contextualização} do próprio aluno como cidadão do mundo ”.
Ressalta -se \textbf{que} o conhecimento geográfico permite “ \textbf{uma} análise crítica da realidade social e natural mais ampla, daí contemplando a diversidade da experiência dos \textbf{homens} na produção do espaço global e dos espaços locais ”, como reforça Cavalcanti ( 2012:37 ).
A proposição curricular da Geografia constante na BNCC já tem sido apresentada em vários estudos, tais como Cavalcanti ( 2002, 2012 ), Callai ( 2000, 2004, 2005 ), Castellar ( 2005 ), Straforini ( 2002, 2004, 2014, 2018 ), e anteriormente por Cholley ( 1942 ) apud Straforini ( 2018:177 ), ao destacar o “ movimento de constituição da Geografia \textbf{enquanto} conhecimento científico \textbf{que} busca, em última \textbf{instância}, desvelar as condições ou as ‘ construções lógicas do presente ’ ”.
Desse modo, a Geografia escolar proposta deve “ possibilitar leituras reflexivas e críticas do mundo, ou ainda ser capaz de formar o cidadão crítico transformador ”, assevera Straforini ( 2018:177 ).
Controvérsias à parte, sobre a construção e o ordenamento legal da BNCC na educação brasileira, os educadores estão sendo provocados para refletir sobre as possibilidades \textbf{que} a Geografia oferece para a formação integral do aluno na Educação Básica.
Ponce ( 2016:1146 ) nos alerta \textbf{que} “ o tempo é, \textbf{hoje}, \textbf{um} bem escasso. \textbf{Um} bem \textbf{que} flui e esgota.
Mas \textbf{um} olhar atento possibilita entendê -lo também como espaço de possibilidades, de oportunidades de \textbf{viver} de modo mais significativo ”.
A ressignificação proposta em \textbf{uma} Geografia \textbf{que} valida o pensamento espacial e o \textbf{raciocínio} geográfico promove \textbf{um} trabalho educativo mais significativo para professores e alunos. ¹ 1 ) Conhecimento ; 2 ) Pensamento científico, crítico e criativo ; 3 ) Repertório cultural ; 4 ) Comunicação ; 5 ) Cultura digital ; 6 ) Trabalho e projeto de vida ; 7 ) Argumentação ; 8 ) Autoconhecimento e autocuidado ; 9 ) Empatia e cooperação ; 10 ) Responsabilidade e cidadania.
Nesse contexto, pensamento espacial e \textbf{raciocínio} geográfico são dois conceitos \textbf{que} podem induzir aprendizagens \textbf{que} devem ser implementadas no desenvolvimento das competências previstas para o ensino de Geografia.
Discutidos primeiramente nos EUA, os documentos orientavam mecanismos para \textbf{que} os alunos pudessem “ aprender a pensar espacialmente ”, a partir de processos cognitivos e sistemáticos de mapeamento com tecnologias da informação geográfica.
A BNCC aponta a \textbf{centralidade} \textbf{que} esses conceitos devem assumir na Geografia escolar, tendo em vista \textbf{que} essas dimensões ampliam \textbf{uma} leitura \textbf{crítico-reflexiva} do mundo, tornando mais efetivas as possibilidades de formar o cidadão transformador, com autonomia em suas relações socioespaciais ( STRAFORINI, 2018 ).
Importante o entendimento da BNCC \textbf{que} recomenda o desenvolvimento do \textbf{raciocínio} geográfico por meio do \textbf{estímulo} ao exercício do pensamento espacial, com a aplicação dos princípios da : analogia, conexão, diferenciação, distribuição, extensão, localização e ordem.
Analogia, quando \textbf{um} fenômeno geográfico sempre é comparável a outros.
A identificação das semelhanças entre fenômenos geográficos é o início da compreensão da unidade terrestre.
Conexão, \textbf{um} fenômeno geográfico nunca acontece isoladamente, mas sempre em interação com outros fenômenos próximos ou distantes.
Diferenciação é a variação dos fenômenos de interesse da Geografia pela superfície terrestre.
Por exemplo, o clima, é resultado das diferenças entre áreas.
Distribuição exprime como os objetos se repartem pelo espaço.
Extensão é o espaço finito e contínuo delimitado pela ocorrência dos fenômenos geográficos.
Localização, posição particular de \textbf{um} objeto na superfície terrestre.
A localização pode ser absoluta ( definida por \textbf{um} sistema de coordenadas geográficas ) ou relativa ( expressa por meio de relações espaciais topológicas ou por interações espaciais ).
Ordem ou arranjo espacial é o princípio geográfico de maior complexidade.
Refere -se ao modo de estruturação do espaço de acordo com as regras da própria sociedade \textbf{que} o produziu ( BRASIL, 2017:358, \textbf{grifos} nossos ).
Acrescenta -se \textbf{que} a intencionalidade do uso dos conceitos pensamento espacial e \textbf{raciocínio} geográfico pode ser identificado, seguindo o \textbf{exposto} por Straforini ( 2018:180 apud De Miguel ( 2016:14 ), a partir de procedimentos voltados para o desenvolvimento de habilidades espaciais, a saber : Visualização espacial – capacidade de manipular, rotacionar, girar ou inverter mentalmente \textbf{estímulos} visuais bi e/tridimensionais ; Orientação espacial – capacidade de imaginar como seria \textbf{um} objeto em \textbf{uma} orientação ou perspectiva diferente do \textbf{sujeito} observador ; Relações espaciais – para o \textbf{autor}, esta é a mais importante das três \textbf{categorias} porque \textbf{implica} a aquisição e o desenvolvimento de processos cognitivos espaciais como, reconhecer as distribuições espaciais, estabelecer associações, identificar padrões de organização e \textbf{hierarquias} no espaço, estabelecer associações e \textbf{correlações} entre fenômenos \textbf{que} têm determinada distribuição espacial ( DE MIGUEL, 2016:14, grifo nosso ).
As habilidades espaciais devem levar em conta a capacidade de autonomia do aluno para articular a subjetividade dos seus sentidos com as atividades de seu cotidiano.
Desenvolver a capacidade relacional de articular, associar, identificar padrões e \textbf{hierarquias} são condições básicas para \textbf{um} aprendizado espacial \textbf{que} vislumbre a \textbf{totalidade}.
Straforini ( 2018 ) ainda contribui para essa compreensão ao destacar \textbf{que} à Geografia interessa a relação sociedade-natureza, como sintetiza Moreira ( 2007:116 ), “ a relação homem-meio é o eixo epistemológico da Geografia ”.
Assim, são as interações sociais, econômicas, humanas, políticas e culturais no espaço, para “ além dos aspectos topológicos ”, \textbf{que} devem ser priorizadas no desenvolvimento do \textbf{raciocínio} geográfico dos alunos.
Nesse percurso, seguindo o entendimento de Freire ( 2005 ), é necessário \textbf{que} haja diálogo e construção de saberes para \textbf{que} os alunos se edifiquem como \textbf{sujeitos} críticos e reflexivos, constituindo -se, acima de tudo, como cidadãos socialmente proativos e responsáveis.
Pretende -se, portanto, a elaboração de \textbf{uma} consciência cidadã sobre o contexto socioespacial do aluno, \textbf{que} tenha como ponto de partida a realidade material e objetiva, bem como as peculiaridades pertinentes ao modo de vida em \textbf{que} está inserido, contribuindo ativamente na compreensão do processo \textbf{dialético} de produção e reprodução do espaço geográfico.
A Geografia ganha \textbf{uma} amplitude em sua \textbf{abordagem} \textbf{conceitual}, procedimental, atitudinal ao oferecer ao aluno \textbf{uma} compreensão de sua atuação no espaço, da participação do(s) outro(s), dos arranjos \textbf{que} se estabelecem nas interações sociais.
Ao desenvolver o pensamento espacial usando o \textbf{raciocínio} geográfico dos alunos, eles efetivam \textbf{uma} leitura do espaço local, a partir da qual podem analisar e questionar a realidade local, regional, global, sendo capazes de observar as diferenciações espaciais, as distribuições dos fenômenos geográficos e suas implicações locais nos mais variados contextos socioespaciais. \textbf{Ademais}, é \textbf{imprescindível} a compreensão de \textbf{que} o conceito de espaço e o de tempo são indissociáveis, \textbf{uma} vez \textbf{que} há a necessidade de serem analisados.
Em se tratando da Geografia, o tempo é \textbf{uma} construção social, \textbf{que} se relaciona às identidades sociais, formadas pelas memórias dos \textbf{sujeitos}.
Analogamente, deve -se utilizar o mesmo princípio para a evolução da natureza : o tempo \textbf{revela} as transformações naturais ocorridas e, por \textbf{conseguinte}, a memória dos acontecimentos \textbf{que} suscitaram no meio físico atualmente \textbf{conhecido}.
Desse modo, a reflexão da temporalidade das ações humanas e naturais representa \textbf{um} desafio no aprendizado da Geografia, a partir relação entre tempo e espaço.
Objetivando o desenvolvimento do \textbf{raciocínio} geográfico, nesse cenário, é importante ressaltar os princípios geográficos \textbf{que} devem ser trabalhados em associação com as principais \textbf{categorias} de análise da Geografia, dos quais é possível destacar : - Espaço geográfico : é \textbf{um} conjunto indissociável de sistemas de objetos ( cidades, vias de transporte, sistemas de produção, etc. ) e sistemas de ações ( fluxos de pessoas, mercadorias, pessoas, práticas políticas, culturais, entre outras ).
Assim, o espaço geográfico constitui -se como o objeto de estudo da Geografia, representando a dimensão de \textbf{totalidade}. - Paisagem : além de referir -se à relação visível do espaço, a paisagem representa as heranças das constantes interações entre \textbf{homem} e natureza, em determinado intervalo de tempo, constituindo -se como \textbf{um} conjunto de objetos concretos.
Nessa perspectiva, essa \textbf{categoria} é transtemporal, juntando objetos passados e presentes, numa construção transversal. - Território : dimensão do espaço \textbf{que} leva em consideração as relações de poder, ou seja, o território é produzido espaço-temporalmente pelas relações de poder engendradas por determinados grupos sociais.
Dessa maneira, pode ser temporário ou permanente e se efetiva em diferentes escalas.
O território leva em conta as dimensões espaciais e seus aspectos econômicos, políticos, culturais e simbólicos. - Lugar : está ligado à construção social \textbf{que} se dá ao longo do tempo, ou seja, \textbf{uma} construção \textbf{sócio-histórica} \textbf{que} se vincula à lógica social.
Ele \textbf{diz} respeito ao espaço onde \textbf{vivemos} e construímos nossa identidade, possibilitando compreender as relações entre o local e o global. - Região : associa -se ao princípio da espacialidade diferencial dos processos sociais e naturais.
Assim, a região se estabelece a partir de \textbf{um} conjunto de elementos políticos, econômicos, culturais e naturais necessários à diferenciação de determinados espaços em diferentes lugares.
Para tanto, é necessário \textbf{que} a aprendizagem da ciência geográfica seja \textbf{vista} não só como a \textbf{mera} capacidade de realizar descrições de informações cotidianas, mas como o domínio de conceitos \textbf{que} permitam a compreensão do mundo de maneira \textbf{holística}, \textbf{despertando} \textbf{um} pensamento crítico \textbf{que} considere as múltiplas realidades existentes.
Além disso, é necessária \textbf{uma} visão multiescalar dos fenômenos de modo \textbf{que} se compreenda desde \textbf{uma} visão local, \textbf{uma} visão global e suas várias conexões espaciais e temporais.

% matched lemmas: abordagem, ademais, assimilação, atual, autor, categoria, centralidade, conceitual, conhecido, conseguinte, contextualização, correlação, crítico-reflexivo, despertar, dialético, dizer, enquanto, estímulo, exposto, grifo, hierarquia, hoje, holístico, homem, implicar, imprescindível, instância, mero, que, raciocínio, revelar, sujeito, sócio-histórico, totalidade, um, ver, vivar, viver
\end{textsample}
